\section{Preliminaries} \label{sec:pre}
In this section, we describe some useful polynomials that will be used to prove our main result. For any integer $k\geq 2$, let $(F_i^{(k)})_{i\geq 0}$ be a sequence of orthogonal polynomials defined by the three-term recurrence relation: 
\begin{equation*}
F_0^{(k)}(x)=1, \qquad F_1^{(k)}(x)=x, \qquad F_2^{(k)}(x)=x^2 - k, 
\end{equation*}
and
\begin{equation} \label{eq:3-term}
F_i^{(k)}(x)=x F_{i-1}^{(k)}(x)- (k-1) F_{i-2}^{(k)}(x)
\end{equation}
for $i\geq 3$. 
The notation $F_i^{(k)}$ is abbreviated to $F_i$ for the rest of the paper. Let $q=\sqrt{k-1}$. The polynomials $(F_i)_{i\geq 0}$ form a sequence of orthogonal polynomials 
with respect to the positive weight
\[
w(x)=\frac{\sqrt{4q^2-x^2}}{k^2-x^2}
\]
on the interval $[-2q, 2q]$ 
(see~\cite[Section~4]{HOb}). 
The polynomials $F_i(q y)/ q^i$ in $y$ are called Geronimus
polynomials~\cite{G30,H15}. 
It follows from~\eqref{eq:3-term} that 
%\begin{align*}
%&F_0^{(k)}(x)=1,& &F_1^{(k)}(x)=x,& 
%&F_2^{(k)}(x)=x^2 - k,& \\
%&F_3^{(k)}(x)=x^3-(2k-1)x,& 
%&F_4^{(k)}(x)=x^4-(3k-2)x^2+k(k-1),&
%&%F_5^{(k)}(x)=x^5 -(4k-3)x^3+(3k-1)(k-1)x, 
%&
%\end{align*}
%and
\begin{equation} \label{eq:3-term2}
F_i(x)=(x^2-2k+2) F_{i-2}(x)- (k-1)^2 F_{i-4}(x)
\end{equation}
for $i\geq 5$. Note that for any $i\geq 0$, $F_{2i}(x)$ and $F_{2i+1}(x)$ are 
even and odd functions of $x$, respectively.

For $i\geq 0$, let $\mathscr{F}_{0,i}(x)=F_{2i}(\sqrt{x})$ and $\mathscr{F}_{1,i}(x)=F_{2i+1}(\sqrt{x})/\sqrt{x}$. 
It follows that $x^{\epsilon} \mathscr{F}_{\epsilon,i}(x^2)=F_{2i+\epsilon}(x)$ for $\epsilon\in \{0,1\}$. 
By~\eqref{eq:3-term2}, the polynomials $\mathscr{F}_{0,i}(x)$ and $\mathscr{F}_{1,i}(x)$ satisfy the following properties:
\begin{gather*}
\mathscr{F}_{0,0}(x)=1, \qquad 
\mathscr{F}_{0,1}(x)=x-k, \qquad
\mathscr{F}_{0,2}(x)=x^2-(3k-2)x+k(k-1), 
\\
\mathscr{F}_{1,0}(x)=1, \qquad 
\mathscr{F}_{1,1}(x)=x-(2k-1),
\end{gather*}
and 
\begin{equation} \label{eq:new_3-term}
\mathscr{F}_{\epsilon, i}(x)=(x-2k+2) \mathscr{F}_{\epsilon,i-1}(x)-(k-1)^2 \mathscr{F}_{\epsilon, i-2}(x)
\end{equation}
for any $i\geq 3$ if $\epsilon=0$, and $i\geq 2$ if $\epsilon=1$. 
Note that $k^{\epsilon} \mathscr{F}_{\epsilon,i}(k^2)=F_{2i+\epsilon}(k)=k(k-1)^{2i-1+\epsilon}=(k-1)^{2i-1+\epsilon}+(k-1)^{2i+\epsilon}$ for $2i+\epsilon \ne 0$. 
For $\epsilon\in \{0,1\}$, the polynomials $(\mathscr{F}_{\epsilon,i})_{i\geq 0}$ form a sequence of orthogonal polynomials with respect to the positive weight 
\[
w_\epsilon(x)=\frac{x^{\epsilon-1/2} \sqrt{4q^2-x}}{k^2-x}
\]
on the interval $[0, 4q^2]$. 

For $i\geq 0$, let $G_i(x)=\sum_{j=0}^{\floor{i/2}}F_{i-2j}(x)$. A simple calculation implies that
\begin{equation} \label{eq:G_i}
{G}_{i}(x)=
\frac{{F}_{i+2}(x)-(k-1)^2 {F}_{i}(x) }{x^2-k^2}
\end{equation}
for $i\geq 1$.
From Lemmas 3.3 and 3.5 in~\cite{CK07}, the polynomials $(G_i)_{i\geq 0}$ form a sequence of 
orthogonal polynomials with respect to 
the positive weight $(k^2-x^2)w(x)=\sqrt{4q^2-x^2}$ on the interval $[-2q,2q]$. From~\eqref{eq:3-term}, we deduce that
\[
G_i(x)=xG_{i-1}(x)-(k-1)G_{i-2}(x)
\]
for $i\geq 2$. 

Let $\mathscr{G}_{\epsilon, i}(x)$ denote the 
polynomial 
\begin{equation} \label{eq:def_scrG}
\mathscr{G}_{\epsilon, i}(x)=\sum_{j=0}^i
\mathscr{F}_{\epsilon, j}(x). 
\end{equation} 
It follows that 
$x^{\epsilon} \mathscr{G}_{\epsilon, j}(x^2)=G_{2j+\epsilon}(x)$. Using~\eqref{eq:G_i}, the polynomial $\mathscr{G}_{\epsilon, i}(x)$ can be expressed as
\begin{equation} \label{eq:pro_scrG}
\mathscr{G}_{\epsilon, i}(x)=
\frac{\mathscr{F}_{\epsilon, i+1}(x)-(k-1)^2\mathscr{F}_{\epsilon, i}(x) }{x-k^2}
\end{equation}
for any $i\geq 2$ if $\epsilon=0$, and
$i\geq 1$ if $\epsilon=1$. 
From Lemmas 3.3 and 3.5 in~\cite{CK07}, 
for $\epsilon\in \{0,1\}$, the polynomials $(\mathscr{G}_{\epsilon,i})_{i\geq 0}$ form a sequence of orthogonal polynomials with respect to 
the positive weight $(k^2-x)w_{\epsilon}(x)=x^{\epsilon-1/2}\sqrt{4q^2-x}$ on the interval $[0,4q^2]$. 

\begin{lemm} \label{lem:coe}
Let $p_l(i,j)$ be the coefficients in 
$x^{\epsilon} \mk \mathscr{F}_{\epsilon,i}(x)\mathscr{F}_{\epsilon,j}(x)\Mk =\Mk \sum_{l=0}^{i+j+\epsilon}\mk p_l(i,j)
\mathscr{F}_{0,l}(x)\mk $. Then
we have 
$p_0(i,j)=
k^{\epsilon}\mathscr{F}_{\epsilon,i}(k^2)\delta_{i,j}$, and 
$p_l(i,j)\geq 0$ for any $l,i,j$. Moreover 
$p_l(i,j)>0$ if and only if $|i-j|\leq l \leq i+j+\epsilon$. 
\end{lemm}
\begin{proof}
We have that
\[
F_{2i+\epsilon}(x)F_{2j+\epsilon}(x)\Mk =\Mk 
x^{2\epsilon}\mathscr{F}_{\epsilon,i}(x^2)\mathscr{F}_{\epsilon,j}(x^2)
\Mk =\Mk 
\sum_{l=0}^{i+j+\epsilon}p_l(i,j)
\mathscr{F}_{0,l}(x^2)
\Mk =\Mk 
\sum_{l=0}^{i+j+\epsilon}p_l(i,j)F_{2l}(x).
\]
By Theorem 3 in~\cite{N}, we obtain that 
$p_0(i,j)=F_{2i+\epsilon}(k) \delta_{i,j}=
k^{\epsilon}\mathscr{F}_{\epsilon,i}(k^2)\delta_{i,j}$, and 
$p_l(i,j)\geq 0$ for any $l,i,j$. Moreover 
$p_l(i,j)>0$ if and only if $|i-j|\leq l \leq i+j+\epsilon$. 
\end{proof}
Let $\Gamma$ be a connected regular bipartite graph. 
The adjacency matrix $\bA$ of $\Gamma$ can be expressed by 
\[
\boldsymbol{A}=
\begin{pmatrix}
\boldsymbol{O} & \boldsymbol{N}\\
\boldsymbol{N}^{\top} & \boldsymbol{O} 
\end{pmatrix}, 
\]
where $\boldsymbol{N}^{\top}$ is the transpose matrix of $\boldsymbol{N}$. 
The matrix $\boldsymbol{N}$ is called the \emph{biadjacency matrix} of $\Gamma$. 
It is not hard to see that 
\begin{equation} \label{eq:AB}
F_{2i}(\boldsymbol{A})= 
\begin{pmatrix}
\mathscr{F}_{0,i}(\boldsymbol{N}\boldsymbol{N}^{\top}) & \boldsymbol{O} \\
\boldsymbol{O} & \mathscr{F}_{0,i}(\boldsymbol{N}^{\top} \boldsymbol{N})
\end{pmatrix}. 
\end{equation}
%If $i$ is odd, then
%\[
%F_i(\boldsymbol{A})= 
%\begin{pmatrix}
%\boldsymbol{O} &\mathscr{F}_i(\boldsymbol{N}\boldsymbol{N}^{\top}) \boldsymbol{N} \\
% \mathscr{F}_i(\boldsymbol{N}^{\top} \boldsymbol{N})\boldsymbol{N}^{\top} & \boldsymbol{O} 
%\end{pmatrix}, 
%\]
%which is needed to write??
Since each entry of $F_{2i}(\boldsymbol{A})$ is non-negative~\cite{S66}, 
each entry of $\mathscr{F}_{0,i}(\boldsymbol{N}\boldsymbol{N}^{\top})$ is also non-negative. 
%From the positivity of $\mathscr{F}_{0,i}(\boldsymbol{N}\boldsymbol{N}^{\top})$, 
%we can obtain the linear programming bound for regular bipartite graphs. 


\section{Linear programming bound for bipartite regular graphs} 
\label{sec:lp}

In this section, we give a linear programming bound for bipartite regular graphs. For general regular graphs, a linear programming bound was obtained by Nozaki~\cite{N}. %As we will see in Table~\ref{tab:1} and Example~\ref{ex:1}, the {\em bipartite} linear programming bound below gives an improvement over the general linear programming bound from~\cite{N}. 
\begin{theo}
\label{thm:lp_bound}
Let $\Gamma$ be a connected bipartite $k$-regular graph with $v$ vertices. 
Let $\{\pm \tau_0, \ldots, \pm \tau_d \}$ be the set of distinct eigenvalues of $\Gamma$, 
where $\tau_0=k$. 
If there exists a polynomial $f(x)=\sum_{i= 0}^t f_i \mathscr{F}_{0,i}(x)$ such that 
$f(k^2)>0$, $f(\tau_i^2) \leq 0$ for each $i\in \{1,\ldots, d \}$, $f_0>0$, and $f_j \geq 0$ for each $j \in \{1,\ldots, t \}$, then 
\begin{equation} 
\label{eq:lp} 
v \leq \frac{2 f(k^2)}{f_0}. 
\end{equation}
Equality holds if and only if for each $i \in \{1,\ldots, d\}$, $f(\tau_i^2)=0$ and for each $j\in \{1,\dots,t\}$, $\tr (f_j\mathscr{F}_{0,j}(\boldsymbol{N}\boldsymbol{N}^{\top}))=0$, 
and $\tr (f_j\mathscr{F}_{0,j}(\boldsymbol{N}^{\top} \boldsymbol{N}))=0$, 
where $\boldsymbol{N}$ is the biadjacency matrix of $\Gamma$.
If equality holds and $f_j>0$ for each $j \in \{1,\ldots, t\}$, then the girth of $\Gamma$ is at least 
$2t+2$. 
\end{theo}
\begin{proof}
From the spectral decomposition $\boldsymbol{N}\boldsymbol{N}^{\top}= 
\sum_{i=0}^d \tau_i^2 \boldsymbol{E}_i$, we deduce that 
\begin{equation} \label{eq:pr1}
f(k^2) \boldsymbol{E}_0+ \sum_{i=1}^d f(\tau_i^2) \boldsymbol{E}_i
=f(\boldsymbol{N}\boldsymbol{N}^{\top})=\sum_{i= 0}^t f_i \mathscr{F}_{0,i}(\boldsymbol{N}\boldsymbol{N}^{\top})=
f_0 \boldsymbol{I} + \sum_{i= 1}^t f_i \mathscr{F}_{0,i}(\boldsymbol{N}\boldsymbol{N}^{\top}), 
\end{equation}
where $\boldsymbol{I}$ is the identity matrix, $\boldsymbol{E}_0=(2/v)\boldsymbol{J}$, and $\boldsymbol{J}$ is the all-ones matrix. 
Taking traces in both sides of~\eqref{eq:pr1}, we get that
\begin{multline*}
 f(k^2)=\tr ( f(k^2) \boldsymbol{E}_0 )
\geq \tr \left( f(k^2) \boldsymbol{E}_0+ \sum_{i=1}^d f(\tau_i^2) \boldsymbol{E}_i\right) \\
= 
\tr \left( f_0 \boldsymbol{I} + \sum_{j= 1}^t f_j\mathscr{F}_{0,j}(\boldsymbol{N}\boldsymbol{N}^{\top}) \right) 
\geq \tr ( f_0 \boldsymbol{I})
=\frac{v f_0}{2}. 
\end{multline*}
Therefore, $v\leq 2 f(k^2)/f_0$. 
By using $\mathscr{F}_{0,j}(\boldsymbol{N}^{\top} \boldsymbol{N})$, 
we can obtain the same bound as~\eqref{eq:lp}. 

If equality holds in~\eqref{eq:lp}, then for each $i \in \{1,\ldots, d\}$, $f(\tau_i^2)=0$ and for each $j\in \{1,\dots,t\}$, $\tr (f_j\mathscr{F}_{0,j}(\boldsymbol{N}\boldsymbol{N}^{\top}))=0$ and $\tr (f_j\mathscr{F}_{0,j}(\boldsymbol{N}^{\top} \boldsymbol{N}))=0$. 
 For the adjacency matrix $\boldsymbol{A}$, the $(u,v)$-entry of $F_{j}(\boldsymbol{A})$ is the number of non-backtracking walks of length $j$ from $u$ to $v$~\cite{S66}. 
Since~\eqref{eq:AB} and $f_j>0$ for each $j \in \{1,\ldots, t\}$, there is no non-backtracking walk of length $2j$ from $u$ to $v$ for each $j \in \{1,\ldots, t\}$. Since $\Gamma$ is bipartite, the girth of $\Gamma$ is at least $2t+2$. 
\end{proof}
%\begin{remark}
%We may normalize $f_0=1$ in Theorem~\ref{thm:lp_bound}. 
%\end{remark}
%\begin{remark} \label{rem:2}

%\end{remark}

\section{Upper bound for bipartite graphs with given second eigenvalue}
\label{sec:newbound}

In this section, we obtain an upper bound on $b(k,\theta)$ using the {\em bipartite} linear programming bound given by Theorem~\ref{thm:lp_bound}. Let $c>0$ be a real number and $t\geq 4$ be an integer. Let $\boldsymbol{B}(k,t,c)$ be the $t \times t$ tridiagonal 
matrix 
with lower diagonal $(1,\ldots,1,c,k)$, upper diagonal $(k,k-1,\ldots, k-1,k-c)$, and constant row sum $k$. 
Let
\[
\boldsymbol{B}(k,3,1)=\begin{pmatrix}
0&k&0 \\
1&0&k-1 \\
0&k& 0
\end{pmatrix}.
\]


